\documentclass[final]{article}

\usepackage{APS360}
\usepackage{amsmath}
\usepackage[utf8]{inputenc} % allow utf-8 input
\usepackage[T1]{fontenc}    % use 8-bit T1 fonts
\usepackage{hyperref}       % hyperlinks
\usepackage{url}            % simple URL typesetting
\usepackage{booktabs}       % professional-quality tables
\usepackage{amsfonts}       % blackboard math symbols
\usepackage{nicefrac}       % compact symbols for 1/2, etc.
\usepackage{microtype}      % microtypography
\usepackage{xcolor}         % colors

\DeclareMathOperator{\sinc}{sinc}

\begin{document}

%==============================================================================
\section{Introduction}
%==============================================================================

X-ray scattering produces a 1-D intensity profile $I(q)$ that encodes the 3-D shape of a
molecule — but computing it exactly via the Debye equation scales as $O(N^2)$ in atom count,
making it prohibitively slow for large structures and incompatible with gradient-based
optimization. This proposal seeks to answer the question 
\emph{can a graph neural network learn to predict $I(q)$ directly from
atomic coordinates?} To answer this, this project aims to train a message-passing GNN as a fast, 
differentiable surrogate for the exact computation. 

%==============================================================================
\section{Background \& Related Work}
%==============================================================================
Graph neural networks for molecular property prediction were pioneered by
SchNet~\cite{schutt_schnet_2018}, which introduced continuous-filter convolutions on
interatomic distances and remains a strong baseline for scalar molecular properties.
DimeNet~\cite{gasteiger_dimenet_2020} extended this by incorporating bond angles, improving
accuracy on quantum-chemical targets such as the QM9 benchmark~\cite{ramakrishnan_qm9_2014}.
On the scattering side, CRYSOL~\cite{svergun_crysol_1995} established the spherical-harmonic
decomposition approach that we use to generate ground-truth $I(q)$ labels.
More recently, \citet{Schutt2021} introduced PaiNN, an equivariant message-passing network
that achieves state-of-the-art accuracy on molecular property benchmarks while remaining
computationally efficient — providing context for the architectural choices made here.

%==============================================================================
\section{Background: X-Ray Scattering and the Debye Equation}
%==============================================================================

When a beam of X-rays is fired at a molecule, each atom deflects (``scatters'') the beam by an
amount that depends on the atom's identity and the deflection angle~$q$. Each atom has a
characteristic \emph{atomic form factor} $f_k(q)$ encoding how strongly it scatters at
angle~$q$. By measuring the total scattered intensity across all angles, we obtain a
\emph{scattering profile} $I(q)$ which is a unique fingerprint of the molecule's 3-D shape.

The exact formula, the \textbf{Debye equation}, sums the pairwise interaction of every atom with
every other atom:
\begin{equation}
    I(q) = \sum_{\alpha}\sum_{\beta}
    f_{\alpha}(q)\,f^{*}_{\beta}(q)\;\sinc\!\bigl(q\,|\mathbf{r}_{\alpha} - \mathbf{r}_{\beta}|\bigr)
    \label{eq:debye}
\end{equation}
where $\sinc(x)=\sin(x)/x$ captures the intuition that atoms far apart interfere less strongly.
Because every pair must be evaluated, this scales as $O(N^2)$ in the number of atoms.
\textbf{CRYSOL}~\cite{svergun_crysol_1995} approximates Eq.~\eqref{eq:debye} by decomposing the
molecular shape into a sum of angular ``modes'' (spherical harmonics), truncated at order~$L$:
\begin{equation}
    I(q) = \sum_{l=0}^{L}\sum_{m=-l}^{l}
    \bigl|A_{lm}(q) - \rho_0\,C_{lm}(q) + \delta\rho\,B_{lm}(q)\bigr|^2
    \label{eq:crysol}
\end{equation}
where $A_{lm}$, $C_{lm}$, $B_{lm}$ encode contributions from the molecule itself, the
displaced solvent, and the surrounding hydration shell, respectively. This reduces the cost to
$O(N \cdot L^2)$, with $L \leq 20$ sufficient for most proteins. A live demo is available at:
(ADD DEMO LINK HERE).

This raises the central question of this project: \emph{can a neural network learn to predict
$I(q)$ directly from the 3-D atomic coordinates of a molecule?}

%==============================================================================
\section{Proposed Architecture}
%==============================================================================

Ground truth intensities are pre-computed via the Stuhrmann decomposition (Eq.~\ref{eq:stuhrmann})
and stored per molecule; $B_{lm}$ coefficients are discarded after computing $I(q)$.

\begin{equation}
    B_{lm}(q) = \sum_i f_i(q)\,j_l(qr_i)\,Y^*_{lm}(\theta_i,\phi_i),
    \qquad
    I(q) = (4\pi)^2 \sum_{l,m} |B_{lm}(q)|^2
    \label{eq:stuhrmann}
\end{equation}

The proposed network takes a molecule's atomic positions and element types as input and
predicts $\hat{I}(q) \in \mathbb{R}^{|Q|}$. Let $\lambda_1$ be the hidden dimension and
$\lambda_2$ the number of message-passing rounds (both treated as tunable hyperparameters).
The network proceeds in four stages:

\begin{enumerate}
    \item \textbf{Embedding.} Each atom's element type is looked up in a trainable table,
    producing a vector $\mathbf{h}_i \in \mathbb{R}^{\lambda_1}$. This gives the network a
    learnable representation of each element.

    \item \textbf{Message passing ($\lambda_2$ rounds).} Each atom updates its vector by
    aggregating information from nearby neighbours. Interatomic distances are encoded as
    Gaussian radial basis functions before being passed to a small MLP (with Swish activations)
    that weights each neighbour's contribution. A residual connection is added at each round.
    After $\lambda_2$ rounds, each atom's vector encodes its local chemical environment.

    \item \textbf{Pooling.} Atom vectors are averaged into a single fixed-size molecule
    vector $\mathbf{h}_{\mathrm{mol}} \in \mathbb{R}^{\lambda_1}$, independent of molecule size.

    \item \textbf{Output head.} A multi-layer MLP with Swish hidden activations maps
    $\mathbf{h}_{\mathrm{mol}}$ to $\hat{I}(q) \in \mathbb{R}^{|Q|}$. A Softplus activation
    on the final layer enforces strict positivity $\hat{I}(q) > 0$, consistent with the
    physical requirement that scattering intensity is always positive.
\end{enumerate}

The network is trained to minimise a log-scale MSE loss,
\begin{equation}
    \mathcal{L} = \frac{1}{|Q|}\sum_{q}
    \bigl(\log\hat{I}(q) - \log I(q)\bigr)^2,
\end{equation}
which accounts for the orders-of-magnitude variation in intensity across the dataset.

A key design consideration is rotational invariance. The $B_{lm}$ coefficients themselves
are rotationally variant; rotating a molecule produces distinct $B_{lm}$ coefficients. 
To predict $B_{lm}$ directly would require rotationally equivariant architectures 
such as e3nn~\cite{e3nn}. However, these models are very computationally expensive and outside of the scope 
of this project. Therefore, it was decided to estimate the magnitudes of the coefficients which are rotationally 
invariant, which has the added bonus of making them real instead of complex.

\section{Data Processing}

A dataset of \textbf{1,139,747 molecules} was assembled from 13 source collections spanning
diverse structural regimes: small clusters (Ar, Ne, Si/Ge, binary metal), crystal structures
(COD, 532,617 entries), small organics (QM9, 133,844), transition-metal complexes
(tmQM, 108,541), biomolecular assemblies (RCSB, 198,173), viral proteins
(viro3D, 85,063), hydration shells (48,571), and metal-organic frameworks (30,871).
In total, 474 unique atom/ion types are represented (charge-resolved, e.g.\ \texttt{O} and
\texttt{O2-} are distinct). Full source attributions are provided in \texttt{sources.tsv},
embedded in the database.

Raw structures were obtained in various formats (.cif, .pdb, .xyz\@) and converted to a
uniform Cartesian XYZ representation. Dummy and non-physical atom sites were removed,
deuterium was mapped to hydrogen, and structures declaring zero atoms were discarded.
Non-element placeholders (\texttt{X}, \texttt{M}, \texttt{Cn}, etc.) were excluded from
the ionic makeup. $I(q)$ ground truths were then computed via the Stuhrmann decomposition
at $L{=}50$ and stored in a compressed HDF5 database.

%==============================================================================
\section{Baseline Model}
%==============================================================================

The natural baseline is the exact Stuhrmann decomposition itself, that is: given atomic coordinates,
$I(q)$ can be computed analytically at $O(N \cdot L^2)$ cost with no approximation error.
This is the same computation used to generate ground-truth labels, so a perfect model would
exactly match it. We report this exact computation as the baseline and measure the GNN's
approximation error relative to it across each dataset regime.

%==============================================================================
\section{Ethical Considerations}
%==============================================================================

All training data is open-source and fully attributed; no private or sensitive information
is involved. The primary risk is misuse of model predictions: since the GNN is an
approximation, predicted $I(q)$ curves could mislead structural biology or materials
research if treated as exact without cross-validation against the ground-truth computation.

\end{document}
